\documentclass[../../../include/open-logic-section]{subfiles}

\begin{document}

\olfileid[id]{sth}{ordinals}{opps}
\olsection{Ordinal Suksesor dan Limit}

Dalam beberapa bab berikutnya, kita akan banyak menggunakan ordinal. Karena
itu, akan berguna jika kita memperkenalkan beberapa gagasan sederhana.

\begin{defn}
Untuk sembarang ordinal $\alpha$, \emph{suksesor}-nya adalah
$\ordsucc{\alpha} =\alpha \cup \{\alpha\}$. Kita mengatakan bahwa $\alpha$
merupakan suatu ordinal \emph{suksesor} jika $\ordsucc{\beta} = \alpha$
untuk suatu ordinal~$\beta$. Kita mengatakan bahwa $\alpha$ merupakan suatu
ordinal \emph{limit} jika dan hanya jika $\alpha$ tidak kosong dan bukan
ordinal suksesor.
\end{defn}
\noindent
Hasil berikut menunjukkan bahwa ini merupakan gagasan \emph{suksesor} yang
tepat:

\begin{prop}
Untuk sembarang ordinal $\alpha$:
\begin{enumerate}
	\item $\alpha \in \ordsucc{\alpha}$;
	\item $\ordsucc{\alpha}$ merupakan suatu ordinal;
	\item tidak ada ordinal $\beta$ sedemikian sehingga $\alpha \in \beta
	\in \ordsucc{\alpha}$.
	\end{enumerate}
\end{prop}

\begin{proof}
Secara trivial, $\alpha \in \alpha \cup \{\alpha\} =
\ordsucc{\alpha}$. Demikian pula, $\ordsucc{\alpha}$ merupakan suatu
himpunan transitif dari ordinal-ordinal, dan karena itu merupakan suatu
ordinal menurut \olref[basic]{corordtransitiveord}. Tidak mungkin pula bahwa
$\alpha \in \beta \in \ordsucc{\alpha}$, sebab dalam hal itu berlaku
$\beta \in \alpha$ atau $\beta = \alpha$, bertentangan dengan
\olref[basic]{ordordered}.
\end{proof}
\noindent
Hal ini juga mengabsahkan suatu varian pembuktian dengan induksi transfinit:

\begin{thm}[Induksi Transfinit Sederhana]\ollabel{simpletransrecursion}
Misalkan $\phi(x)$ suatu formula sedemikian sehingga:
\begin{enumerate}
	\item $\phi(\emptyset)$; dan
	\item untuk sembarang ordinal $\alpha$, jika $\phi(\alpha)$ maka
	$\phi(\ordsucc{\alpha})$; dan
	\item jika $\alpha$ merupakan suatu ordinal limit dan $(\forall \beta \in
	\alpha)\phi(\beta)$, maka $\phi(\alpha)$.
\end{enumerate}
Maka $\forall \alpha \phi(\alpha)$.
\end{thm}

\begin{proof}
Kita membuktikan kontraposisinya. Jadi, andaikan terdapat suatu ordinal yang
memenuhi $\lnot\phi$; misalkan $\gamma$ merupakan ordinal terkecil semacam
itu. Maka berlaku salah satu dari: $\gamma = \emptyset$; atau $\gamma =
\ordsucc{\alpha}$ untuk suatu $\alpha$ sedemikian sehingga $\phi(\alpha)$;
atau $\gamma$ merupakan suatu ordinal limit dan $(\forall \beta \in
\gamma)\phi(\beta)$.
\end{proof}
\noindent
Sedikit notasi terakhir akan berguna kemudian:

\begin{defn}\ollabel{defsupstrict}
Jika $X$ merupakan suatu himpunan ordinal, maka $\supstrict(X) =
\bigcup_{\alpha \in X} \ordsucc{\alpha}$.
\end{defn}
\noindent
Di sini, ``lsub'' merupakan singkatan dari ``least strict upper bound''
(batas atas ketat terkecil).\footnote{Beberapa buku menggunakan
``$\text{sup}(X)$'' untuk gagasan ini. Namun, buku lain menggunakan
``$\text{sup}(X)$'' untuk batas atas \emph{tak ketat} terkecil, yakni cukup
$\bigcup X$. Jika $X$ mempunyai anggota terbesar, $\alpha$, kedua
gagasan ini berbeda: batas atas \emph{ketat} terkecil adalah
$\ordsucc{\alpha}$, sedangkan batas atas \emph{tak ketat} terkecil hanyalah
$\alpha$.} Hasil berikut menjelaskan hal ini:

\begin{prop}
Jika $X$ merupakan suatu himpunan ordinal, $\supstrict(X)$ merupakan ordinal
terkecil yang lebih besar daripada setiap ordinal dalam $X$.
\end{prop}

\begin{proof}
Misalkan $Y = \Setabs{\ordsucc{\alpha}}{\alpha \in X}$, sehingga
$\supstrict(X) = \bigcup Y$. Karena ordinal bersifat transitif dan setiap
anggota suatu ordinal merupakan suatu ordinal, $\supstrict(X)$ merupakan
suatu himpunan transitif dari ordinal-ordinal, dan karena itu merupakan suatu
ordinal menurut \olref[basic]{corordtransitiveord}.

Jika $\alpha \in X$, maka $\ordsucc{\alpha} \in Y$, sehingga
$\ordsucc{\alpha} \subseteq \bigcup Y = \supstrict(X)$, dan karena itu
$\alpha \in \supstrict(X)$. Jadi, $\supstrict(X)$ lebih besar secara ketat
daripada setiap ordinal dalam $X$.

Sebaliknya, jika $\alpha \in \supstrict(X)$, maka $\alpha \in
\ordsucc{\beta} \in Y$ untuk suatu $\beta \in X$, sehingga $\alpha \leq
\beta \in X$. Jadi, $\supstrict(X)$ merupakan batas atas ketat
\emph{terkecil} bagi $X$.
\end{proof}

\end{document}
